\documentclass[11pt,a4paper]{article}

\usepackage{fontspec}
\setmainfont{DejaVu Serif}[
  BoldFont={DejaVu Serif Bold},
  ItalicFont={DejaVu Serif},
  ItalicFeatures={FakeSlant=0.18},
  BoldItalicFont={DejaVu Serif Bold},
  BoldItalicFeatures={FakeSlant=0.18}
]
\setsansfont{DejaVu Sans}
\usepackage{polyglossia}
\setmainlanguage{english}
\usepackage[a4paper,margin=20.5mm]{geometry}
\usepackage{microtype}
\usepackage{setspace}
\setstretch{1.035}
\usepackage{amsmath,amssymb,amsthm,mathtools,bm}
\usepackage{booktabs,longtable,array,tabularx}
\usepackage{enumitem}
\usepackage{xcolor}
\usepackage{tikz}
\usetikzlibrary{arrows.meta,positioning,calc}
\usepackage[unicode,hidelinks]{hyperref}
\usepackage[nameinlink,noabbrev]{cleveref}
\emergencystretch=3em
\allowdisplaybreaks

\hypersetup{
  pdftitle={SI-HEP Quantity Passport for Relativistic Beam Observation},
  pdfauthor={Stas, Independent Research Program},
  pdfsubject={A reversible dimensional contract between SI, HEP natural units, and accelerator mixed units},
  pdfkeywords={natural units, SI, HEP, dimensional analysis, LHC, magnetic rigidity, four-momentum}
}

\definecolor{obsblue}{RGB}{37,83,138}
\definecolor{obsred}{RGB}{150,45,45}
\definecolor{obsgreen}{RGB}{30,115,78}
\definecolor{lightblue}{RGB}{239,246,253}
\definecolor{lightred}{RGB}{253,242,242}
\definecolor{lightgreen}{RGB}{239,249,244}

\theoremstyle{definition}
\newtheorem{definition}{Definition}[section]
\theoremstyle{plain}
\newtheorem{theorem}[definition]{Theorem}
\newtheorem{proposition}[definition]{Proposition}
\newtheorem{corollary}[definition]{Corollary}
\theoremstyle{remark}
\newtheorem{remark}[definition]{Remark}
\newtheorem{warning}[definition]{Boundary}

\newcommand{\R}{\mathbb{R}}
\newcommand{\dd}{\mathop{}\!d}
\newcommand{\abs}[1]{\left\lvert #1\right\rvert}
\newcommand{\norm}[1]{\left\lVert #1\right\rVert}
\newcommand{\Ezero}{\mathcal E_0}
\newcommand{\USI}{\mathcal U_{\mathrm{SI}}}
\newcommand{\Nhep}{\mathcal N_{\Ezero}}
\newcommand{\GeV}{\mathrm{GeV}}
\newcommand{\MeV}{\mathrm{MeV}}
\newcommand{\SI}{\mathrm{SI}}
\newcommand{\HEP}{\mathrm{HEP}}
\newcolumntype{L}[1]{>{\raggedright\arraybackslash}p{#1}}
\newcolumntype{Y}{>{\raggedright\arraybackslash}X}

\title{\bfseries SI--HEP Quantity Passport\\
\large A reversible dimensional contract for relativistic beam observation}
\author{Stas\\
\small Independent Research Program; ORCID: 0009-0000-2294-705X}
\date{GO Core extension v0.5 --- 28 July 2026}

\begin{document}
\maketitle

\begin{center}
\fcolorbox{obsblue}{lightblue}{%
\parbox{0.92\linewidth}{\small
\textbf{Status.} Normative infrastructure for GO Core v0.2 and the LHC
beam module. The passport fixes the meaning of \(c=\hbar=1\), gives an
invertible SI--HEP chart on the mechanical sector, and defines exact
accelerator-unit bridges. It does not propose a new physical law.}}
\end{center}

\begin{abstract}
Natural units suppress conversion constants in formulas but do not
erase physical dimensions. For a mechanical quantity with SI dimension
\(L^aM^bT^\tau\), fix an energy anchor \(\Ezero>0\), set
\[
 d_E=b-a-\tau,\qquad u=a+\tau,\qquad v=a-2b,
\]
and define the SI representative of one natural unit by
\[
 \USI(a,b,\tau;\Ezero)=\hbar^u c^v\Ezero^{d_E}.
\]
Then
\(\Nhep(Q)=Q/\USI(a,b,\tau;\Ezero)\) is dimensionless and the inverse
map is multiplication by the same representative. This gives a
multiplicative, exponent-additive, reversible chart on the \(LMT\)
sector. It does not automatically cover current, temperature, amount
of substance, luminous intensity, or electromagnetic field
normalization. For LHC beam kinematics the missing electromagnetic
step is bypassed by the exact mixed-unit identity
\[
 p[\GeV/c]=0.299792458\,z\,B\rho[\mathrm{T\,m}],
\qquad q=ze.
\]
The passport also separates the SI momentum chart
\(p^\mu=(E/c,\bm p)\), the SI energy chart
\(\Pi^\mu=(E,c\bm p)\), and the HEP numerical chart. Consequently the
usual shorthand \(E^2=\bm p^2+m^2\) is admitted only after every
component has been represented in the same energy unit. Exact
definition-based conversions, measured constants, and operational
machine parameters are assigned distinct uncertainty statuses.
\end{abstract}

\tableofcontents

\section{Scope and failure mode}

The blocked LHC manuscript used four locally familiar notations:
\[
 E\,[\mathrm{TeV}],\qquad
 p\,[\GeV/c],\qquad
 B\rho\,[\mathrm{T\,m}],\qquad
 E^2=\bm p^2+m^2\quad(c=1).
\]
Each can be correct. Their untyped composition is not a quantity
calculus. In particular, the last equation does not permit an SI mass
in kilograms to be added to a numerical momentum in \(\GeV/c\).

\begin{warning}[Claim firewall]
Setting \(c=\hbar=1\) is a coordinate choice on the dimension group,
not an empirical statement that the constants have ceased to exist.
It neither fixes an electromagnetic rationalization convention nor
makes tesla an energy squared without further definitions.
\end{warning}

The present passport covers:
\begin{enumerate}[label=(\roman*)]
  \item the mechanical \(LMT\) sector;
  \item relativistic four-momentum in a declared metric signature;
  \item charged-particle magnetic rigidity through a mixed-unit bridge;
  \item RF voltage-to-energy conversion through a declared charge state;
  \item uncertainty and provenance fields needed for reproducibility.
\end{enumerate}

\section{Typed quantity records}

\begin{definition}[Quantity record]
A quantity record is
\[
\mathsf q=
(\mathsf{id},Q,\bm d,\mathsf{kind},\mathsf{frame},
 \mathsf{context},\mathsf{representation},\mathsf U),
\]
where \(Q\) is the physical quantity, \(\bm d\) is its seven-component
SI dimension vector, \(\mathsf{kind}\) is its semantic type,
\(\mathsf{frame}\) is a reference-frame label. The context obeys
\[
\mathsf{context}\in
\{\SI,\ \HEP\text{-natural},\ \text{accelerator-mixed}\}.
\]
The field \(\mathsf{representation}\) specifies which powers of
\(c\) and \(\hbar\) have been absorbed, and
\(\mathsf U\) contains exactness, uncertainty, and provenance.
\end{definition}

Equality and addition are permitted only for records with compatible
physical dimension, semantic type, frame law, and representation.
Multiplication combines dimension vectors and representations.
A context change is an explicit map; it is never inferred from a unit
label alone.

\begin{center}
\begin{tikzpicture}[
  node distance=10mm and 14mm,
  box/.style={draw,rounded corners,minimum width=39mm,minimum height=12mm,align=center},
  arr/.style={-{Latex[length=2mm]},thick}
]
\node[box,fill=lightblue] (si) {SI quantity\\\(Q,\ L^aM^bT^\tau\)};
\node[box,fill=lightgreen,right=of si] (hep) {HEP coordinate\\\(\Nhep(Q),\ E^{d_E}\)};
\node[box,fill=lightred,below=of $(si)!0.5!(hep)$] (pass) {passport\\\(\hbar^u c^v\Ezero^{d_E}\)};
\draw[arr] (si) -- node[above]{naturalize} (hep);
\draw[arr] (hep) -- node[below]{restore} (si);
\draw[arr] (pass) -- (si);
\draw[arr] (pass) -- (hep);
\end{tikzpicture}
\end{center}

\section{Mechanical naturalization theorem}

Let the admissible dimension lattice be
\[
\mathcal D_{\mathrm{mech}}
=\{(a,b,\tau,0,0,0,0):a,b,\tau\in\mathbb Q\}.
\]
The last four exponents must vanish. Fractional exponents are allowed
when a positive branch is part of the quantity definition.

\begin{theorem}[Unique \(c,\hbar\) reduction to an energy dimension]
For \(D=L^aM^bT^\tau\), there is a unique triple
\((d_E,u,v)\in\mathbb Q^3\) such that
\[
 D=[\hbar]^u[c]^v[E]^{d_E}.
\]
It is
\[
 \boxed{d_E=b-a-\tau,\qquad u=a+\tau,\qquad v=a-2b.}
\]
\end{theorem}

\begin{proof}
Using
\([E]=ML^2T^{-2}\), \([\hbar]=ML^2T^{-1}\), and
\([c]=LT^{-1}\), equality of the \(M,L,T\) exponents gives
\[
 d_E+u=b,\qquad
 2d_E+2u+v=a,\qquad
 -2d_E-u-v=\tau.
\]
This linear system has determinant \(1\), hence a unique solution.
Solving yields the displayed exponents.
\end{proof}

\begin{definition}[Natural coordinate and inverse]
Fix an energy anchor \(\Ezero\), here \(1\,\GeV\). Define
\[
\USI(D;\Ezero)=
\hbar^{a+\tau}c^{a-2b}\Ezero^{\,b-a-\tau},
\qquad
\Nhep(Q)=\frac{Q}{\USI(D;\Ezero)}.
\]
The inverse is
\[
Q=\Nhep(Q)\,\USI(D;\Ezero).
\]
\end{definition}

\begin{proposition}[Algebraic consistency]
For admissible positive quantities \(Q_1,Q_2\) and rational \(r\),
\[
\Nhep(Q_1Q_2)=\Nhep(Q_1)\Nhep(Q_2),\qquad
\Nhep(Q_1^r)=\Nhep(Q_1)^r.
\]
For quantities of the same complete type,
\(\Nhep(Q_1+Q_2)=\Nhep(Q_1)+\Nhep(Q_2)\).
\end{proposition}

\begin{proof}
The exponent maps \(D\mapsto(d_E,u,v)\) are linear, so the SI
representatives multiply under dimension addition. Division by the
representative proves the product and power laws. The addition law
holds only when both terms share the same representative.
\end{proof}

If the anchor changes from \(\Ezero\) to
\(\Ezero'=\lambda\Ezero\), then
\[
\mathcal N_{\Ezero'}(Q)=\lambda^{-d_E}\mathcal N_{\Ezero}(Q).
\]
The physical quantity is unchanged; its numerical coordinate is
covariant with its energy dimension.

\section{Standard conversion rows}

The defining values
\[
c=299\,792\,458\ \mathrm{m\,s^{-1}},\quad
h=6.62607015\times10^{-34}\ \mathrm{J\,s},\quad
e=1.602176634\times10^{-19}\ \mathrm C
\]
are exact in the SI. Hence
\[
1\,\mathrm{eV}=e(1\,\mathrm V)
=1.602176634\times10^{-19}\,\mathrm J,
\qquad
1\,\GeV=1.602176634\times10^{-10}\,\mathrm J
\]
are exact.
The displayed decimal for \(\hbar=h/(2\pi)\) is rounded, although
\(\hbar\) is an exactly defined quantity.

\begin{center}
\small
\begin{tabularx}{\linewidth}{@{}L{31mm}L{29mm}L{43mm}Y@{}}
\toprule
Quantity & HEP unit & SI representative & Exactness \\
\midrule
energy & \(\mathrm{GeV}\) & \(1.60217663400000\times 10^{-10}\ \mathrm{J}\) & definition-exact; decimal rounded \\
mass & \(\mathrm{GeV}/c^2\) & \(1.78266192162790\times 10^{-27}\ \mathrm{kg}\) & definition-exact; decimal rounded \\
momentum & \(\mathrm{GeV}/c\) & \(5.34428599267831\times 10^{-19}\ \mathrm{kg\,m\,s^{-1}}\) & definition-exact; decimal rounded \\
length & \(\mathrm{GeV}^{-1}\) & \(1.97326980459302\times 10^{-16}\ \mathrm{m}\) & definition-exact; decimal rounded \\
duration & \(\mathrm{GeV}^{-1}\) & \(6.58211956950907\times 10^{-25}\ \mathrm{s}\) & definition-exact; decimal rounded \\
area / cross section & \(\mathrm{GeV}^{-2}\) & \(3.89379372171859\times 10^{-32}\ \mathrm{m^2}\) & definition-exact; decimal rounded \\
\bottomrule
\end{tabularx}

\end{center}

Useful equivalent forms are
\[
\hbar c=197.326980459302\ \MeV\,\mathrm{fm},
\qquad
1\,\GeV^{-2}=0.389379372172\ \mathrm{mb}.
\]
The time row uses angular frequency:
\[
E=\hbar\omega,\qquad
1\,\GeV\longleftrightarrow
1.519267447878626\times10^{24}\ \mathrm{rad\,s^{-1}}.
\]
For ordinary frequency one must instead use
\(E=h\nu\), giving
\[
1\,\GeV\longleftrightarrow
2.417989242084918\times10^{23}\ \mathrm{Hz}.
\]
The factor \(2\pi\) is part of the passport and cannot be discarded.

\section{Four-momentum and invariant-mass passports}

Adopt metric signature \(\eta=\operatorname{diag}(1,-1,-1,-1)\).
There are three representations of the same four-momentum.

\begin{longtable}{@{}L{36mm}L{55mm}L{55mm}@{}}
\toprule
Representation & Definition & Mass shell \\
\midrule
\endhead
SI momentum chart &
\(p^\mu_{\SI}=(E/c,\bm p)\) &
\(p_{\SI,\mu}p^\mu_{\SI}=m^2c^2\) \\
SI energy chart &
\(\Pi^\mu_{\SI}=c\,p^\mu_{\SI}=(E,c\bm p)\) &
\(\Pi_{\SI,\mu}\Pi^\mu_{\SI}=(mc^2)^2\) \\
HEP numerical chart &
\(P^\mu_{\HEP}=\Pi^\mu_{\SI}/\Ezero\) &
\(P_{\HEP,\mu}P^\mu_{\HEP}=(mc^2/\Ezero)^2\) \\
\bottomrule
\end{longtable}

Thus the collider shorthand
\[
m^2=E^2-\norm{\bm p}^2
\]
is legal only when \(m\) means the energy-valued mass \(mc^2\),
\(\bm p\) means \(c\bm p_{\SI}\), and all three numerical entries use
the same energy anchor. The passport-preferred notation is
\[
\mu:=mc^2,\qquad \bm\pi:=c\bm p_{\SI},\qquad
\mu^2=E^2-\norm{\bm\pi}^2.
\]

For a reconstructed system,
\[
\mu_{\mathrm{inv}}^2
=
\left(\sum_i E_i\right)^2
-\left\lVert\sum_i\bm\pi_i\right\rVert^2.
\]
Equivalently in SI momentum coordinates,
\[
m_{\mathrm{inv}}^2c^4
=
\left(\sum_i E_i\right)^2
-c^2\left\lVert\sum_i\bm p_i\right\rVert^2.
\]
These are the same equation under the explicit map
\(\bm\pi_i=c\bm p_i\).

\section{Exact accelerator bridges}

\subsection{Magnetic rigidity}

For transverse motion in a magnetic field,
\[
\abs{q}\,B\rho=\abs{\bm p}_{\SI}.
\]
Let \(q=ze\), with signed dimensionless charge state \(z\).
Write \(e_0=e/(1\,\mathrm C)\) and
\(c_0=c/(1\,\mathrm{m\,s^{-1}})\) for the exact numerical values in
the declared SI units. Then
\[
q=z e_0\,\mathrm C,\qquad
1\,\GeV=10^9e_0\,\mathrm J,\qquad
1\,\GeV/c=\frac{10^9e_0}{c_0}\ \mathrm{kg\,m\,s^{-1}},
\]
division by the momentum unit gives
\[
\boxed{
p[\GeV/c]=0.299792458\,z\,B\rho[\mathrm{T\,m}].
}
\]
The coefficient is exact for the declared units because
\[
0.299792458=\frac{c[\mathrm{m\,s^{-1}}]}{10^9}.
\]
Signs encode bending orientation; magnitudes use \(\abs z\).

\begin{warning}[Geometry boundary]
The dipole bending radius is not \(C/(2\pi)\) for a ring with straight
sections. A circumference may determine the revolution frequency, but
it cannot be substituted for \(\rho\) in \(B\rho\) without a lattice
model.
\end{warning}

\subsection{RF voltage}

For a particle crossing an idealized RF gap at phase \(\phi\),
\[
\Delta E[\mathrm J]=q[\mathrm C]\,V[\mathrm V]\sin\phi.
\]
With \(q=ze\), the exact unit identity is
\[
\boxed{
\Delta E[\mathrm{eV}]=z\,V[\mathrm V]\sin\phi,\qquad
\Delta E[\GeV]=10^{-9}z\,V[\mathrm V]\sin\phi.
}
\]
The conversion is exact. The single-gap sinusoidal model is an
approximation to a real cavity and has a distinct model-error status.

\section{Electromagnetic normalization boundary}

The mechanical theorem requires zero exponents of \(I,\Theta,N,J\).
It therefore does not naturalize coulombs, amperes, tesla, or volts
by itself. In a rationalized Heaviside--Lorentz field convention one
often writes
\[
e_{\mathrm{HL}}=\sqrt{4\pi\alpha},
\]
whereas other conventions distribute factors of \(4\pi\) differently
between fields and couplings. The statements
\[
c=\hbar=1,\qquad B\sim E^2
\]
do not select among these conventions.

The normative LHC policy is therefore:
\begin{enumerate}[label=(\alph*)]
  \item retain the Lorentz force and \(B\) in SI;
  \item encode charge as \(q=ze\);
  \item cross to HEP momentum only through the rigidity bridge;
  \item do not register a standalone tesla-to-\(\GeV^2\) conversion.
\end{enumerate}
This is a deliberate restriction, not a missing calculation.

\section{Exactness and uncertainty}

\begin{longtable}{@{}L{38mm}L{50mm}L{59mm}@{}}
\toprule
Status & Examples & Rule \\
\midrule
\endhead
Definition-exact &
\(c,h,e,\mathrm{eV}\leftrightarrow\mathrm J\),
rigidity coefficient &
Adds no physical uncertainty; printed decimals may be rounded. \\
Measured constant &
proton mass energy &
Retain the source value, standard uncertainty, edition, and covariance
when relevant. \\
Operational reference &
LHC beam energy, bunch population &
Record epoch, machine mode, nominal/measured status, and precision. \\
Rounded geometry &
public circumference \(26659\,\mathrm m\) &
Do not report derived revolution frequency as more accurate than the
input geometry. \\
Model-dependent &
single-gap RF gain, ideal linear optics &
Separate model discrepancy from unit conversion. \\
\bottomrule
\end{longtable}

For a differentiable derived vector \(\bm y=f(\bm x)\), the first-order
covariance rule is
\[
\Sigma_y=J_f\Sigma_xJ_f^{\mathsf T}.
\]
Exact conversion constants enter \(J_f\) but contribute no independent
variance. Correlated inputs must not be propagated as independent
errors.

\section{Normative machine contract}

A valid conversion event must contain:
\[
\begin{aligned}
(&\mathsf{quantity\_id},\mathsf{semantic\_kind},
\bm d_{\SI},\mathsf{frame},\\
 &\mathsf{source\_context},\mathsf{source\_unit},
\mathsf{target\_context},\mathsf{target\_unit},\\
 &\mathsf{absorbed\_constants},\Ezero,
\mathsf{representation},\mathsf{exactness},
\mathsf{uncertainty},\mathsf{provenance}).
\end{aligned}
\]
The conversion is rejected if:
\begin{itemize}
  \item a residual \(I,\Theta,N,J\) dimension crosses into the mechanical
  HEP chart without an extension contract;
  \item \(m\), \(p\), and \(E\) are combined without a common
  four-momentum representation;
  \item a field is converted without an electromagnetic convention;
  \item an operational LHC value lacks an epoch or mode;
  \item a rounded public input is marked exact.
\end{itemize}

\section{Claim audit}

\begin{center}
\begin{tabularx}{\linewidth}{@{}L{47mm}L{28mm}Y@{}}
\toprule
Statement & Status & Reason \\
\midrule
Mechanical exponent map & theorem & linear dimension algebra \\
Round-trip conversion & theorem & inverse by construction \\
Rigidity coefficient \(0.299792458\) & exact identity & SI definitions and declared units \\
\(\mu^2=E^2-\bm\pi^2\) & standard result & Lorentz metric and representation choice \\
Run 3 \(6.8\,\mathrm{TeV}\) reference & operational datum & CERN epoch-dependent value \\
Completeness for all SI quantities & rejected & electromagnetic and other residual dimensions excluded \\
New physical prediction & not claimed & passport is infrastructure \\
\bottomrule
\end{tabularx}
\end{center}

\section*{References}
\addcontentsline{toc}{section}{References}

\begin{enumerate}[label={[\arabic*]}]
  \item BIPM, \emph{The International System of Units (SI)},
  9th ed., version 4.01.
  \url{https://www.bipm.org/documents/20126/41483022/SI-Brochure-9-EN.pdf}
  \item P.\ J.\ Mohr et al., ``CODATA recommended values of the
  fundamental physical constants: 2022,'' 2025.
  \url{https://physics.nist.gov/cuu/pdf/JPCRD2022CODATA.pdf}
  \item Particle Data Group, ``Physical Constants,'' in
  \emph{Review of Particle Physics 2024}.
  \url{https://pdg.lbl.gov/2024/reviews/rpp2024-rev-phys-constants.pdf}
  \item CERN, ``Large Hadron Collider: facts and figures.''
  \url{https://home.cern/science/accelerators/large-hadron-collider/}
  \item CERN, \emph{LHC Design Report}, Vol.\ I, 2004.
  \url{https://cds.cern.ch/record/782076/files/CERN-2004-003-V1-ft.pdf}
\end{enumerate}

\end{document}
